\begin{proofbox}
	\textbf{\textcolor{CProof}{证明从哪里入手？}}

	\textbf{先把对数还原成乘方。} 设 $\log_a M=x$、$\log_a N=y$，就表示
	\[
		a^x=M,\qquad a^y=N.
	\]
	接下来只用已经学过的指数运算法则。

	\medskip
	\textbf{\textcolor{CProof}{逐步证明}}
	\begin{description}[style=nextline, leftmargin=2.8em, labelsep=0.5em, itemsep=0.55em]
		\item[\textbf{第一步（积变加）：}]
			因为 $a^x\cdot a^y=a^{x+y}=MN$，所以使 $a$ 的某次方等于 $MN$ 的指数就是 $x+y$。
			\[
				\log_a(MN)=x+y=\log_a M+\log_a N.
			\]
			\textbf{关键：} 同底幂相乘，指数相加。例：$2^3\cdot2^2=2^5$。

		\item[\textbf{第二步（商变减）：}]
			因为 $a^x\div a^y=a^{x-y}=M/N$，所以
			\[
				\log_a\frac{M}{N}=x-y=\log_a M-\log_a N.
			\]
			\textbf{关键：} 同底幂相除，指数相减。例：$2^3\div2^2=2^1$。

		\item[\textbf{第三步（指数提前）：}]
			因为 $(a^x)^m=a^{mx}=M^m$，所以
			\[
				\log_a M^m=mx=m\log_a M.
			\]
			\textbf{关键：} 幂再乘方，指数相乘。例：$(2^3)^2=2^6$。

		\item[\textbf{第四步（换底）：}]
			设 $\log_a b=t$，即 $a^t=b$。等式两边取以 $c$ 为底的对数，再用第三步：
			\[
				\log_c(a^t)=\log_c b
				\quad\Longrightarrow\quad t\log_c a=\log_c b.
			\]
			因为 $a\neq1$，所以 $\log_c a\neq0$，两边除以 $\log_c a$：
			\[
				\boxed{\log_a b=t=\frac{\log_c b}{\log_c a}}.
			\]
			\textbf{关键：} 只是把原来求的指数 $t$，改用另一个底数算出来。
	\end{description}

	\medskip
	\textbf{\textcolor{CProof}{为什么这套证明值得记？}}

	不必分别死记四条推导：\textbf{对数变乘方 $\longrightarrow$ 用指数法则 $\longrightarrow$ 再写回对数。} 所有步骤都从同一个想法出发。
\end{proofbox}
